% ============================================================================
%  1. MOTIVACIÓN  (0--20 min)
% ============================================================================
\section{Motivación}

% S02
\begin{frame}[label=pregunta]{La productividad sube 0.72\% por un trimestre}
  \textbf{¿Qué hace la economía?}\par
  \vspace{0.45cm}
  ¿Trabajan más o menos los hogares?\par
  \vspace{0.15cm}
  ¿Consumen la ganancia o la ahorran?\par
  \vspace{0.15cm}
  ¿Cuánto dura el efecto?\par
  \vspace{0.6cm}
  Al final de la sesión, cada respuesta será un número, y cada número saldrá de una
  ecuación que ya conocen y de un archivo que habrán escrito ustedes.
  \note{Pedir predicciones a mano alzada: ¿las horas suben o bajan? Anotar la votación
  en la pizarra y volver a ella en el Acto 1 (slide acto1-horas). El 0.72\% es una
  desviación estándar del choque de PTF en la calibración de King y Rebelo (1999).}
\end{frame}

% S03
\begin{frame}[label=pizarra]{Hasta dónde llega la pizarra}
  \begin{columns}[T,onlytextwidth]
    \begin{column}{0.47\textwidth}
      \textbf{Depreciación total, $\delta=1$}\par
      \vspace{0.25cm}
      Solución cerrada (PD06, parte E).\par
      \vspace{0.2cm}
      $\hat h_t=0$: las horas no se mueven.\par
      \vspace{0.2cm}
      $\hat c_t=\hat y_t$: el hogar consume siempre la misma fracción.
    \end{column}
    \begin{column}{0.47\textwidth}
      \textbf{Caso realista, $\delta=0.025$}\par
      \vspace{0.25cm}
      No hay solución cerrada (PD06, E.3).\par
      \vspace{0.2cm}
      Las horas \emph{sí} responden.\par
      \vspace{0.2cm}
      Hay que resolver numéricamente.
    \end{column}
  \end{columns}
  \vspace{0.6cm}
  \begin{block}{La paradoja}
    El mecanismo que hace interesante al RBC aparece justo donde la pizarra deja de
    alcanzar.
  \end{block}
  \note{Conectar con la recitación: los alumnos ya probaron en E.3 que con $\delta<1$ el
  supuesto $\hat h_t=0$ es incompatible con la Euler. La sesión de hoy es la parte G:
  ``beyond the blackboard''.}
\end{frame}

% S04
\begin{frame}[label=laboratorio]{El modelo como laboratorio}
  \begin{quote}
    \enquote{Our task, as I see it \ldots\ is to write a FORTRAN program that will accept
    specific economic policy rules as `input' and will generate as `output' statistics
    describing the operating characteristics of time series we care about.}
  \end{quote}
  \begin{flushright}\small\parencite[pp.~709--710]{lucas1980}\end{flushright}
  \vspace{0.1cm}
  \begin{quote}
    \enquote{The key tool macroeconomists use is the computational experiment.}
  \end{quote}
  \begin{flushright}\small\parencite[p.~341]{kydland2004}\end{flushright}
  \vspace{0.1cm}

  En 1980 el programa de investigación ya se pensaba como código.
  \note{Subrayar la fecha. Kydland y Prescott (1982) convirtieron la visión de Lucas en
  método: calibrar el modelo como un instrumento de medición y luego interrogarlo.}
\end{frame}

% S05
\begin{frame}[label=linea-tiempo]{De la visión a la herramienta}
  \centering
  \begin{tikzpicture}[font=\footnotesize]
    \draw[very thick, upblue] (0,0) -- (9.8,0);
    \foreach \x/\y/\a in {
        0.6/1976/{Crítica\\de Lucas},
        2.3/1980/{Lucas,\\Blanchard--\\Kahn},
        4.0/1982/{Kydland--\\Prescott},
        5.7/1983/{Fair--Taylor},
        7.4/1994/{Juillard crea\\\textbf{Dynare}},
        9.1/2025/{30 años}}{
      \filldraw[upblue] (\x,0) circle (2.4pt);
      \node[upblue, font=\footnotesize\bfseries] at (\x,0.38) {\y};
      \node[align=center, text width=1.6cm, font=\scriptsize] at (\x,-0.8) {\a};
    }
  \end{tikzpicture}
  \par\vspace{0.6cm}
  \raggedright
  \textbf{Blanchard y Kahn (1980)}: cuándo un modelo lineal tiene solución única y
  estable. Es lo que Dynare comprueba con \cod{check}.\par
  \vspace{0.25cm}
  \textbf{Fair y Taylor (1983)}: el método para modelos no lineales de previsión
  perfecta.\par
  \vspace{0.25cm}
  \textbf{Dynare (1994, CEPREMAP)}: Michel Juillard une todo en un solo programa.
  \fuente{\parencite{bk1980,ft1983,kp1982,dynare}}
  \note{Cinco minutos como máximo. En febrero de 2025 el Banque de France celebró los
  treinta años de Dynare con una conferencia.}
\end{frame}

% S06
\begin{frame}[label=que-es-dynare]{¿Qué es Dynare?}
  \begin{bluebox}{Dynare}
    Un \textbf{preprocesador} y una biblioteca de rutinas para MATLAB u Octave que
    resuelve, simula y estima modelos DSGE.
  \end{bluebox}
  \vspace{0.2cm}
  Escribimos el modelo casi como en el papel. Dynare se encarga de:
  \begin{itemize}
    \item Estado estacionario y su verificación.
    \item Linealización con derivadas exactas.
    \item Condición de Blanchard--Kahn y regla de decisión.
    \item Funciones de impulso-respuesta y momentos.
    \item Estimación bayesiana (solo un adelanto al final).
  \end{itemize}
  \fuente{\parencite{dynare}}
  \note{Mencionar alternativas (Python, Julia) y por qué usamos Dynare: sintaxis
  declarativa, estándar en bancos centrales y licencia de MATLAB en la UP.}
\end{frame}

% S07
\begin{frame}[label=objetivos]{Al terminar la sesión}
  \begin{enumerate}
    \item Escribirás un \texttt{.mod} completo del modelo RBC.\par\vspace{0.15cm}
    \item Leerás lo que imprime Dynare: estado estacionario, valores propios y regla de
      decisión.\par\vspace{0.15cm}
    \item Explicarás cada impulso-respuesta con una ecuación.\par\vspace{0.15cm}
    \item Confrontarás el modelo con los datos.
  \end{enumerate}
  \vspace{0.4cm}
  \textbf{Material}\par
  \vspace{0.1cm}
  Handout \emph{Un crash course de Dynare}\par
  Taller \texttt{rbc\_taller.mod}\par
  PD06, parte G
  \note{Comprobar que todos tienen MATLAB abierto y la carpeta del taller descargada.
  Si alguien no tiene Dynare instalado, que trabaje en pareja.}
\end{frame}
